% ===========================================================================
% Séquence — contrat, signature par lien à jeton, activation de l'employé.
%
% Tracé sur backend/app/services/contract_service.py (2026-09-22) :
%   create / send  → token = secrets.token_urlsafe(32), statut SENT
%   get_by_token   → le JETON tient lieu d'authentification
%   sign           → garde d'état, empreinte SHA-256 des termes, certificat,
%                    piste d'audit (nom, date UTC, IP, agent), puis création
%                    idempotente de l'employé (dédoublonnage par e-mail).
%
% Frontière de confiance à montrer : le candidat signe SANS COMPTE. Le jeton
% est le seul justificatif — d'où les gardes d'état explicites avant toute
% écriture.
% ===========================================================================
\begin{tikzpicture}[x=1mm, y=-1mm]

  \def\lifelen{124}

  % ─── Participants ───────────────────────────────────────────────────────
  \node[pqhead] (rh)  at (0,0)   {RH};
  \node[pqhead] (api) at (32,0)  {API contrats};
  \node[pqhead] (sv)  at (66,0)  {Service contrat};
  \node[pqhead] (db)  at (100,0) {PostgreSQL};
  \node[pqhead, draw=piqbitWarn, fill=piqbitWarnLite, text=piqbitWarn]
                (cd)  at (136,0) {Candidat\\{\scriptsize non authentifié}};

  \foreach \n in {rh,api,sv,db}
    \draw[pqlifeline] (\n.south) -- ++(0,\lifelen);
  \draw[pqlifeline, draw=piqbitWarn] (cd.south) -- ++(0,\lifelen);

  % ─── Phase 1 : émission ─────────────────────────────────────────────────
  \draw[pqmsg] (0,14) -- node[pqmsglbl] {rédige et envoie le contrat} (32,14);
  \fill[pqactivation] (30.5,14) rectangle (33.5,40);

  \draw[pqmsg] (33.5,21) -- node[pqmsglbl] {génère un jeton} (66,21);
  \node[pqnote, anchor=north west, text width=34mm] at (66,24)
    {\texttt{token\_urlsafe(32)}\\statut \textsc{sent}};

  \draw[pqmsg] (66,34) -- node[pqmsglbl] {enregistre} (100,34);

  \draw[pqmsg, draw=piqbitWarn] (33.5,45) --
    node[pqmsglbl, text=piqbitWarn] {lien de signature (courriel)} (136,45);

  % ─── Séparateur ─────────────────────────────────────────────────────────
  \draw[draw=piqbitWarn, dash pattern=on 2pt off 2pt] (-8,54) -- (146,54);
  \node[anchor=west, font=\scriptsize\itshape, text=piqbitWarn,
        fill=white, inner sep=2pt] at (-8,54)
    {frontière de confiance — le jeton tient lieu d'authentification};

  % ─── Phase 2 : consultation et signature ────────────────────────────────
  \draw[pqmsg, draw=piqbitWarn] (136,60) --
    node[pqmsglbl, text=piqbitWarn] {ouvre le lien} (33.5,60);
  \fill[pqactivation] (64.5,60) rectangle (67.5,112);

  \draw[pqmsg] (33.5,67) -- node[pqmsglbl] {résout le jeton} (66,67);
  \draw[pqreturn] (66,73) -- (136,73);
  \node[pqnote, anchor=north east, text width=40mm, align=right] at (134,76)
    {termes en lecture seule};

  \draw[pqmsg, draw=piqbitWarn] (136,84) --
    node[pqmsglbl, text=piqbitWarn] {signe (nom + tracé manuscrit)} (67.5,84);

  % ─── Phase 3 : écriture gardée ──────────────────────────────────────────
  % Largeur bornée à 50 mm : au-delà, la boîte mord sur la ligne de vie du
  % service contrat (x = 66) et sur les flèches d'écriture.
  \node[pqwarn, anchor=north west, text width=50mm] at (-2,90)
    {\textbf{Gardes d'état avant toute écriture.}
     Déjà signé $\rightarrow$ \texttt{409}.
     Refusé ou statut autre que \textsc{sent} $\rightarrow$ \texttt{410}.
     Un lien rejoué ne signe pas deux fois.};

  \draw[pqmsg] (66,97) -- node[pqmsglbl] {empreinte \textsc{sha-256} des termes} (100,97);
  \draw[pqmsg] (66,104) -- node[pqmsglbl] {certificat + piste d'audit} (100,104);
  \node[pqnote, anchor=north west, text width=36mm] at (102,107)
    {nom, date UTC, IP, agent};

  \draw[pqmsg] (66,116) -- node[pqmsglbl] {crée l'employé si absent} (100,116);
  \node[pqnote, anchor=north west, text width=36mm] at (102,119)
    {dédoublonnage par e-mail\\ — opération idempotente};

\end{tikzpicture}
